% Źródła: Hartshorne s. 312--317 (§34, równoległość graniczna), 388 nn. (§41, końce).

\section{Promienie równoległe graniczne i końce} % lang-pl
\section{Semirette parallele limite e estremi} % lang-it
\label{sekcja_promienie_graniczne}

Równoległość graniczna była dla Łobaczewskiego, Bolyaia i Gaussa podstawowym narzędziem, ale ani oni, ani Eves nie dowodzą wszystkiego, co potrzeba. % lang-pl
W tej sekcji zobaczymy, że równoległość graniczna jest relacją równoważności, i jak z jej klas zrodzi się pojęcie końca prostej -- punktu idealnego. % lang-pl
Il parallelismo limite era lo strumento fondamentale di Lobachevskij, Bolyai e Gauss, ma né loro né Eves dimostrano tutto ciò che serve. % lang-it
In questa sezione vedremo che il parallelismo limite è una relazione di equivalenza e come dalle sue classi nascerà la nozione di estremo di una retta -- il punto ideale. % lang-it

Hartshorne oznacza półprostą symbolem $Aa$, gdzie $A$ jest jej początkiem, a $a$ oznacza prostą ją zawierającą wraz z wyborem kierunku; dwie półproste są \emph{współkońcowe} (ang.~\emph{coterminal}), jeśli leżą na jednej prostej i „idą w tę samą stronę'' \cite[s. 312]{hartshorne_2000}. % lang-pl
Hartshorne indica una semiretta con il simbolo $Aa$, dove $A$ è la sua origine e $a$ la retta che la contiene con la scelta di un verso; due semirette sono \emph{coterminali} se stanno su una stessa retta e «vanno nella stessa direzione» \cite[p. 312]{hartshorne_2000}. % lang-it

\begin{definition}
[półproste równoległe graniczne] % lang-pl
[semirette parallele limite] % lang-it
\index{półproste!równoległe graniczne} % lang-pl
\index{semirette!parallele limite} % lang-it
	Półprosta $Aa$ jest równoległa graniczna do półprostej $Bb$, jeśli są współkońcowe albo leżą na różnych prostych różnych od $AB$, nie przecinają się i każda półprosta wewnątrz kąta $BAa$ przecina półprostą $Bb$. % lang-pl
	Una semiretta $Aa$ è parallela limite a una semiretta $Bb$ se sono coterminali oppure giacciono su rette distinte diverse da $AB$, non si incontrano e ogni semiretta interna all'angolo $BAa$ interseca la semiretta $Bb$. % lang-it
\end{definition}

Hartshorne \cite[s. 312]{hartshorne_2000} (notacja $Aa \parallel\!\!\parallel Bb$). % lang-pl
Hartshorne \cite[p. 312]{hartshorne_2000} (notazione $Aa \parallel\!\!\parallel Bb$). % lang-it

Istnienie takich półprostych nie wynika z aksjomatów neutralnych; zakłada je dopiero aksjomat Hilberta z sekcji \ref{sekcja_hiperboliczna} \cite[s. 313]{hartshorne_2000}. % lang-pl
L'esistenza di tali semirette non discende dagli assiomi neutrali; la postula solo l'assioma di Hilbert della sezione \ref{sekcja_hiperboliczna} \cite[p. 313]{hartshorne_2000}. % lang-it

\begin{proposition}
	Równoległość graniczna półprostych jest przenoszalna na półproste współkońcowe, symetryczna i przechodnia. % lang-pl
	Il parallelismo limite tra semirette è trasmissibile a semirette coterminali, simmetrico e transitivo. % lang-it
\end{proposition}

Dowody: Hartshorne \cite[s. 313--316]{hartshorne_2000} (prop. 34.9--34.11 i lemat 34.12); u Evesa te same fakty są w sekcji \ref{sekcja_hiperboliczna}. % lang-pl
Dimostrazioni: Hartshorne \cite[p. 313--316]{hartshorne_2000} (prop. 34.9--34.11 e lemma 34.12). % lang-it

\begin{definition}
[koniec] % lang-pl
[estremo] % lang-it
\index{koniec prostej} % lang-pl
\index{estremo di una retta} % lang-it
	Klasę równoważności półprostych równoległych granicznych nazywamy \emph{końcem}. % lang-pl
	Chiamiamo \emph{estremo} una classe di equivalenza di semirette parallele limite. % lang-it
\end{definition}

Hartshorne \cite[s. 316]{hartshorne_2000} (wniosek 34.13). % lang-pl
Hartshorne \cite[p. 316]{hartshorne_2000} (corollario 34.13). % lang-it

Końce to dokładnie punkty idealne z sekcji \ref{sekcja_punkty_idealne}; Hilbert wyprowadzi z nich całe ciało liczbowe, a modele Kleina i Poincarego okażą się jedynymi płaszczyznami hiperbolicznymi (zobacz sekcję \ref{sekcja_ciala_koncow}). % lang-pl
Gli estremi sono esattamente i punti ideali della sezione \ref{sekcja_punkty_idealne}; Hilbert ne ricaverà un intero corpo numerico, e i modelli di Klein e di Poincaré si riveleranno gli unici piani iperbolici (si veda la sezione \ref{sekcja_ciala_koncow}). % lang-it

Zadania: konstrukcje trójkąta granicznego i jego „dwusiecznych'' -- Hartshorne \cite[s. 317]{hartshorne_2000} (zad. 34.9--34.12). % lang-pl
Esercizi: costruzioni del triangolo limite e delle sue «bisettrici» -- Hartshorne \cite[p. 317]{hartshorne_2000} (es. 34.9--34.12). % lang-it
